\documentclass[11pt]{article}
\usepackage[margin=1in]{geometry}
\usepackage{amsmath,amssymb,amsthm}
\usepackage{graphicx}
\usepackage{booktabs}
\usepackage{hyperref}
\usepackage{natbib}
\usepackage{xcolor}
\usepackage{microtype}
\usepackage{multirow}

\newcommand{\Gr}{\mathrm{Gr}}
\newcommand{\BN}{\mathrm{BN}}
\newcommand{\R}{\mathbb{R}}
\newcommand{\TODO}[1]{\textcolor{red}{[TODO: #1]}}
\newcommand{\CITE}[1]{\textcolor{red}{[CITE: #1]}}
\newcommand{\dg}{d_{\mathrm{g}}}
\newcommand{\effdim}{\mathrm{dim}_{\mathrm{eff}}}

\title{A Three-Rung Evaluation of RNA Structure Awareness\\
  in Foundation Models}

\author{Elliot Tower\\
  University of Edinburgh\\
  \texttt{elliot@elliottower.ai}}

\begin{document}
\maketitle

\begin{abstract}
We present a three-rung evaluation framework for measuring RNA
secondary structure awareness in foundation models, each rung demanding
a strictly stronger form of structural knowledge.  Rung~1 tests
stem-versus-loop discrimination under a nucleotide-stratified
permutation null that absorbs GC-enrichment confounds, reducing the
apparent structure signal from ${\sim}5.5\times$ to $1.8\times$ across
52 RNA families.  Rung~2 applies a dinucleotide-stratified null; a
single family (Hepatitis~C IRES domain~III) survives both orders of
composition control.  Rung~3 asks whether a model resolves which
specific position pairs with which: after mutating one base to its
Watson--Crick complement, does the model perturb the paired partner
more than adjacent stem positions?

Among ten models spanning 1.2M--7B parameters and five architectures,
ERNIE-RNA (86M, RNA-pretrained) is the sole model to clear Rung~3,
with mean perturbation specificity $\mathrm{PS} = 0.117$, 28 of 30
qualifying families exceeding a within-stem derangement null, and
partner-is-max precision of 88.3\%.  All other models produce
$\mathrm{PS} < 0.005$.  Randomized-weight ERNIE-RNA yields
$\mathrm{PS} = 2.5 \times 10^{-8}$, confirming the result is learned.
ERNIE-RNA's architectural base-pairing attention bias provides the
capacity; pretraining fills it with specific pairing knowledge.

A case study on HTT mRNA CAG-repeat alleles (six models, five clinical
variants) shows RNA-FM embeddings collapse across repeat counts
(cosine distance $= 0.000$), while ERNIE-RNA tracks the structural
transition monotonically.  All evaluations are preregistered with
SHA-frozen protocols.  The three-rung framework provides a reusable
template for evaluating structure awareness in sequence foundation
models.
\end{abstract}


%% ────────────────────────────────────────────
\section{Introduction}
\label{sec:intro}

Foundation models pretrained on nucleotide sequences have been
proposed for RNA biology, but evaluating whether they encode secondary
structure---the pattern of intramolecular base pairing that governs
RNA stability and function---requires metrics that distinguish learned
representations from sequence-composition artifacts.

We initially adopted the Grassmannian geodesic distance ($\dg$), a
geometric metric for comparing embedding subspaces between wild-type
and mutant sequences.  A zero-parameter 3-mer frequency baseline
matched the trained models' $\dg$ scores, falsifying the metric: it
captured $k$-mer composition, a property that random linear maps
preserve under the Johnson--Lindenstrauss lemma
\citep{johnson1984extensions}, rather than learned structure.  Similar
failure modes have been reported for geometric transfer metrics in
other domains \citep{tower2026boundary}.

The $\dg$ collapse motivated a replacement metric with explicit
composition controls.  We designed a mutation sensitivity ratio under
complement swap (A$\leftrightarrow$U, C$\leftrightarrow$G) with a
nucleotide-stratified permutation null.  The key observation: stems
are enriched for G and C (68.5\% vs 43.9\% in loops), and complement
swaps at G/C positions produce inherently larger embedding
perturbations than at A/U positions regardless of structural context.
Any naive stem-vs-loop comparison inherits this 24.7 percentage-point
enrichment as a confound.

These observations motivated a graded evaluation framework.
``Structure awareness'' encompasses at least three levels of
resolution, each demanding strictly stronger knowledge.  The first
level---stem-versus-loop discrimination---can be achieved by any
representation sensitive to nucleotide composition.  The second
level---discrimination surviving composition controls at one and two
orders of sequence context---requires genuine sensitivity beyond GC
enrichment.  The third level---knowing which specific position pairs
with which---requires that the model encode the pairing topology.

We apply this three-rung ladder to ten models across five
architectures and three orders of magnitude in parameter count.  The
evaluation reveals a surprising asymmetry: the hardest test produces
the clearest positive result.  ERNIE-RNA (86M), the only model with
an architectural base-pairing attention bias, clears Rung~3 with
partner specificity three orders of magnitude above all other models.

Structure probing of RNA models has previously focused on
downstream task performance \citep{rnafm,dalla2023nucleotide} or
attention visualization rather than composition-controlled metrics.
\citet{ethayarajh2019} showed that BERT embeddings become increasingly
anisotropic through layers, a phenomenon we observe in RNA-FM on
repeat sequences.  Vig et al.\ demonstrated that protein language
model attention heads correlate with contact maps
\citep{vig2020bertology}; our attention-contact analysis extends this
to RNA, finding strong learned correlation for NT~v2 and near-zero
for RNA-FM.  The mutation sensitivity approach relates to \emph{in
silico} mutagenesis for variant effect prediction
\citep{frazer2021disease}, applied here to the representation layer
with structure-defined mutation categories (stem vs.\ loop).
Composition confounds in genomic analyses have a long history:
dinucleotide frequencies affect codon usage statistics
\citep{karlin1998dinucleotide}, CpG content confounds methylation
studies, and GC content correlates with gene density and chromatin
accessibility.  Our stratified permutation framework adapts these
insights to the evaluation of neural sequence models.


%% ────────────────────────────────────────────
\section{Methods}
\label{sec:methods}

\subsection{RNA families and models}
\label{sec:families_models}

Phase~1 evaluated 12 RNA families with experimentally determined
secondary structures from PDB crystal structures and Rfam consensus
annotations (Table~\ref{tab:rna_families}).  Phase~2 expands to
$N = 52$ families drawn from Rfam~14.10, spanning ribozymes,
riboswitches, snRNAs, cis-regulatory elements, miRNA precursors,
CRISPR repeats, and other ncRNAs (Appendix~A).  Positions are labeled
stem (base-paired) or loop (unpaired) from dot-bracket annotation.
Minimum criteria: $\geq 20$~nt, $\geq 5$ stem and 5 loop positions.

\begin{table}[t]
\centering
\caption{Phase~1 RNA families.  Source: PDB crystal structures and
  Rfam consensus annotations.  Stem: number of base-paired positions.}
\label{tab:rna_families}
\small
\begin{tabular}{@{}llccl@{}}
\toprule
\textbf{RNA family} & \textbf{Class} & \textbf{Length} &
  \textbf{Stem} & \textbf{Source} \\
\midrule
tRNA$^{\mathrm{Phe}}$ (yeast)  & tRNA         & 76  & 42 & PDB 1EHZ \\
tRNA$^{\mathrm{Ala}}$ (human)  & tRNA         & 76  & 42 & Rfam RF00005 \\
5S rRNA (\textit{E.\ coli})    & rRNA         & 120 & 72 & PDB 1C2X \\
Hammerhead ribozyme             & Ribozyme     & 48  & 16 & PDB 2OEU \\
SAM riboswitch                  & Riboswitch   & 94  & 52 & PDB 2GIS \\
TPP riboswitch                  & Riboswitch   & 80  & 40 & PDB 2CKY \\
SRP RNA helix~8                 & Structural   & 77  & 44 & PDB 1HQ1 \\
miR-21 precursor                & Regulatory   & 72  & 64 & Rfam RF00001 \\
IRES HCV domain~II              & Regulatory   & 77  & 50 & PDB 1P5O \\
HDV ribozyme                    & Ribozyme     & 76  & 48 & PDB 1DRZ \\
RNase~P specificity             & Structural   & 78  & 42 & PDB 1NBS \\
U2 snRNA stem                   & Spliceosomal & 76  & 46 & PDB 3K1V \\
\bottomrule
\end{tabular}
\end{table}

Phases~1--2 evaluate five models spanning attention-based, SSM-based,
and hybrid architectures (Table~\ref{tab:models}): RNA-FM (BERT, 99M,
character, RNA-pretrained; \citealp{rnafm}), NT~v2 (ESM+GLU, 56M,
6-mer, DNA; \citealp{dalla2023nucleotide}), HyenaDNA (Hyena SSM,
5.4M, character, DNA; \citealp{nguyen2024hyenadna}), Caduceus
(BiMamba, 14M, character, DNA; \citealp{schiff2024caduceus}), and Evo
(StripedHyena, 7B, byte-level, DNA; \citealp{nguyen2024evo}).
Phase~3 adds RiNALMo (BERT+RoPE+SwiGLU, 650M, character, RNA;
\citealp{penic2024rinalmo}) and UTR-LM (transformer encoder, 1.2M,
character, UTR-focused; \citealp{chu2024utrlm}).
Phase~4 adds ERNIE-RNA (BERT, 86M, character, RNA;
\citealp{wang2024ernierna}), SpliceBERT (BERT, 19.4M, character, RNA;
\citealp{chen2024splicebert}), and DNABERT-2 (MosaicBERT+ALiBi, 117M,
BPE, DNA; \citealp{zhou2024dnabert2}).  Six of the ten models are
RNA-pretrained; parameter counts span four orders of magnitude.

\begin{table}[t]
\centering
\caption{Models under test.  Phases~1--2 evaluate the first five;
  Phase~3 adds RiNALMo and UTR-LM; Phase~4 adds the final three.}
\label{tab:models}
\small
\begin{tabular}{@{}llcll@{}}
\toprule
\textbf{Model} & \textbf{Architecture} & \textbf{Params} &
  \textbf{Tokenization} & \textbf{Pretraining} \\
\midrule
RNA-FM     & BERT encoder          & 99M  & Character  & RNA \\
NT~v2      & ESM + GLU attention   & 56M  & 6-mer      & DNA \\
HyenaDNA   & Hyena SSM             & 5.4M & Character  & DNA \\
Caduceus   & BiMamba SSM           & 14M  & Character  & DNA \\
Evo        & StripedHyena          & 7B   & Byte-level & DNA \\
\midrule
RiNALMo    & BERT + RoPE + SwiGLU  & 650M & Character  & RNA \\
UTR-LM     & Transformer encoder   & 1.2M & Character  & UTR (RNA) \\
\midrule
ERNIE-RNA  & BERT encoder          & 86M  & Character  & RNA \\
SpliceBERT & BERT encoder          & 19.4M & Character & RNA \\
DNABERT-2  & MosaicBERT + ALiBi    & 117M & BPE        & DNA \\
\bottomrule
\end{tabular}
\end{table}

\subsection{Structure sensitivity metrics}
\label{sec:metrics}

\textbf{Mutation sensitivity ratio.}  For each position in each RNA
sequence, we perform a complement swap (A$\leftrightarrow$U,
C$\leftrightarrow$G) and measure the cosine distance between
wild-type and mutant embeddings at the mutated position.  The
structure sensitivity ratio $R$ is the mean cosine distance at stem
positions divided by the mean at loop positions.  We compute $R$ at
every layer and report the maximum; the null models below are computed
under the same max-over-layers selection.

\textbf{Attention-contact correlation.}  For models with accessible
attention matrices (RNA-FM, NT~v2, RiNALMo, UTR-LM, ERNIE-RNA,
SpliceBERT, DNABERT-2), we compute the Spearman rank correlation
between symmetrized attention weights and a binary contact matrix (1
for base-paired positions, 0 otherwise).  For NT~v2's 6-mer and
DNABERT-2's BPE tokenization, contacts are aggregated to the token
level.  All models except DNABERT-2 are loaded with
\texttt{attn\_implementation=``eager''} to ensure attention weights
are returned.  DNABERT-2's MosaicBERT architecture uses a fused Wqkv
projection; we extract content-only attention by hooking the Wqkv
output, splitting Q and K, and computing
$\mathrm{softmax}(QK^T/\sqrt{d_k})$ without ALiBi position biases.
We compare trained attention to an untrained baseline with
randomized weights.

\textbf{Structure probing.}  Logistic regression on per-position
embeddings predicting stem vs loop, using GroupKFold cross-validation
by RNA family to test cross-family generalization.  Balanced accuracy
is reported against the majority baseline.

\subsection{Composition null models}
\label{sec:nulls}

The nucleotide-stratified null shuffles stem/loop labels independently
within each nucleotide type (A, U, G, C), preserving per-nucleotide
composition while destroying the structure-label association.  For
each of 100 permutations, a single shuffled label assignment is
generated and $R$ is computed at every layer; the maximum across
layers is taken.  The 95th percentile of these 100 max-over-layers
values defines the null threshold.  This controls for first-order
composition (per-nucleotide GC enrichment) but not for dinucleotide
context---stem-G occurs predominantly in GC base pairs, while loop-G
occurs in diverse dinucleotide contexts.

The dinucleotide-stratified null extends this to second-order context:
labels are shuffled within each dinucleotide type (previous nucleotide
+ current nucleotide: AA, AC, AG, \ldots).  The same max-over-layers
procedure applies.  This null is run on all families exceeding the
first-order null, plus tRNA-Ala and SAM riboswitch (Phase~1 survivors)
regardless of Phase~2 result.

Both phases were preregistered with frozen code and analysis SHAs
(Phase~1: \texttt{694b43b}; Phase~2: \texttt{bd4b3fd}).  Phase~1
($N = 12$) tested whether RNA-FM's trained/untrained ratio exceeds
2.0 (H1), whether $\geq 6/12$ families exceed the null (H1b), and
whether probing exceeds the majority baseline (H5).  Phase~2
($N \geq 50$) replicated these at scale (H6, H6b, H8), added a
dinucleotide null test (H7: $\geq 1$ family survives), a confirmatory
sign test promoting the Phase~1 exploratory unanimity (H10: trained
$>$ untrained in $\geq 75\%$ of families), and an NT~v2 attention
test (H11: trained $\rho > 0.15$, untrained $< 0.05$).  A
methodological correction applied to both phases: the Phase~1 null
had been computed per-layer independently, while the metric used
max-over-layers.  Phase~2 corrects this mismatch.


\subsection{Perturbation specificity (Rung~3)}
\label{sec:rung3_methods}

Rung~3 tests whether a model resolves which specific position pairs
with which.  For each Watson--Crick base pair $(i, j)$ in a stem with
$\geq 3$ eligible pairs, we mutate position~$i$ to its complement and
measure the cosine distance $\Delta_k$ between wild-type and mutant
embeddings at every other stem position~$k$.  Perturbation specificity
is
\[
  \mathrm{PS}(i, j) = \Delta_j - \max(\Delta_{j-1},\, \Delta_{j+1}),
\]
where $\Delta_j$ is the distance at the paired partner and
$\Delta_{j \pm 1}$ are distances at adjacent stem positions.  A
positive PS means the model perturbs the correct partner more than its
neighbors.  We compute PS at every layer and report the maximum.

\textbf{Positive control gate.}  Each family must show significantly
greater stem perturbation than loop perturbation (paired $t$-test,
$p < 0.05$).  Families failing this gate are excluded: a model
insensitive to complement swaps in stems cannot be expected to resolve
individual partners.

\textbf{Within-stem derangement null.}  For each stem with
$K \geq 3$ eligible pairs, 1000 derangements permute partner labels
within the stem.  The real PS is compared against the distribution of
deranged PS values.  A family exceeds the null if the observed mean PS
exceeds the 95th percentile of derangement means.

\textbf{H3: Partner-is-max.}  For deep-interior pairs (not at stem
termini), the fraction of pairs where the true partner receives the
maximum perturbation among all stem positions.  The chance baseline is
$1/3$ (partner vs two neighbors); the preregistered threshold is
binomial $> 1/3$.

Phase~6 was preregistered with SHA \texttt{c19aa59} (Bonferroni
$\alpha = 0.05/3 = 0.0167$).  Preregistered hypotheses: H1$_6$
(at least one model with mean PS~$> 0$ and $\geq 7$ families exceeding
the derangement null), H2$_6$ (RNA-pretrained models higher PS than
DNA-pretrained, rank-biserial $> 0.5$), H3$_6$ (at least one model
with partner-is-max fraction exceeding $1/3$).


%% ────────────────────────────────────────────
\section{Results}
\label{sec:results}

\subsection{Attention-contact correlation is architectural}
\label{sec:attn_results}

NT~v2 trained attention correlates with base-pairing contacts at mean
Spearman $\rho = 0.322$ across 52 families
(Table~\ref{tab:attention}).  An earlier analysis using the default
SDPA attention backend reported $\rho = 0.000$ for untrained weights,
but SDPA silently returns \texttt{None} for attention weights,
producing a spurious zero.  With eager attention computation,
untrained NT~v2 produces $\rho = 0.328$---slightly \emph{higher} than
trained.  The per-family delta (trained $-$ untrained) averages
$-0.006$ with no systematic direction: 22 families show higher trained
correlation, 30 show higher untrained.  The attention-contact signal
is an architectural property of the ESM attention mechanism and
6-mer tokenization, present from random initialization.

\begin{table}[t]
\centering
\caption{Attention-contact Spearman correlation ($N = 52$).  NT~v2's
  high correlation is entirely architectural (untrained $\geq$
  trained).  DNABERT-2 shows learned content attention but excludes
  ALiBi position biases ($^*$).}
\label{tab:attention}
\begin{tabular}{@{}lccl@{}}
\toprule
\textbf{Model} & \textbf{Trained} & \textbf{Untrained} &
  \textbf{Interpretation} \\
\midrule
NT~v2      & 0.322 & 0.328 & Architectural \\
DNABERT-2$^*$  & 0.257 & 0.000 & Learned (content-only) \\
RiNALMo    & 0.101 & 0.082 & Weak ($\Delta = +0.019$) \\
ERNIE-RNA  & 0.099 & 0.075 & Weak ($\Delta = +0.024$) \\
SpliceBERT & 0.064 & 0.056 & None ($\Delta = +0.008$) \\
RNA-FM     & 0.051 & 0.060 & None \\
UTR-LM     & 0.052 & 0.053 & None \\
\bottomrule
\end{tabular}
\end{table}

Among character-tokenized models, RiNALMo and ERNIE-RNA show the
strongest suggestions of learned attention: RiNALMo trained $\rho =
0.101$ versus untrained $0.082$ ($\Delta = +0.019$), ERNIE-RNA
trained $\rho = 0.099$ versus untrained $0.075$ ($\Delta = +0.024$).
Both deltas are small.  SpliceBERT (0.064 vs 0.056), RNA-FM (0.051
vs 0.060), and UTR-LM (0.052 vs 0.053) show no learned signal.
DNABERT-2's content-only attention (computed from Wqkv projections
without ALiBi position biases; Section~\ref{sec:phase4}) shows the
largest trained/untrained gap ($0.257$ vs $0.000$), but this metric
excludes the position-dependent biases that dominate the model's
actual attention computation.

NT~v2's 6-mer tokenization explains the architectural correlation.
Base-pairing partners---typically 3--30 nucleotides apart---are
aggregated into adjacent tokens, producing structured attention
patterns from positional proximity alone.  Random weights with this
tokenization already capture enough pairwise structure to match
trained attention.

\subsection{Scale does not compensate for domain mismatch}
\label{sec:scale}

Evo (7B parameters, byte-level tokenization, DNA-pretrained) produces
the weakest mutation sensitivity of all trained models
(Table~\ref{tab:phase1_summary}).  Two families (tRNA-Phe: 0.927,
tRNA-Ala: 0.899) show ratios below 1.0---Evo is less sensitive to stem
mutations than loop mutations for these structures.  Probing accuracy
(0.615) peaks at layer~0 (embedding) and falls below Caduceus (0.628,
14M) and HyenaDNA (0.621, 5.4M).  Despite being $500\times$ larger
than Caduceus and $70\times$ larger than RNA-FM, Evo shows no
advantage on any structure metric.  This observation is consistent
across mutation sensitivity, probing, and the null comparison, though
$n = 5$ models is too few to support a general scaling claim.

\begin{table}[t]
\centering
\caption{Mutation sensitivity summary, Phase~1 ($N = 12$).}
\label{tab:phase1_summary}
\begin{tabular}{@{}lccc@{}}
\toprule
\textbf{Model} & \textbf{Mean ratio} &
  \textbf{Families $>$ null} & \textbf{Probing $\Delta$} \\
\midrule
RNA-FM (99M, RNA)    & 1.774 & 1/12 & +0.013 \\
Caduceus (14M, DNA)  & 1.257 & 3/12 & +0.041 \\
HyenaDNA (5.4M, DNA) & 1.223 & 0/12 & +0.034 \\
Evo (7B, DNA)        & 1.170 & 0/12 & +0.028 \\
NT~v2 (56M, DNA)     & 1.094 & 6/12 & +0.016 \\
\bottomrule
\end{tabular}
\end{table}

\subsection{Composition controls absorb the embedding-level signal}
\label{sec:composition}

The nucleotide-stratified null reduced the apparent RNA-FM signal from
${\sim}5.5\times$ (naive, before composition control) to $1.8\times$
(composition-controlled).  The 24.7 percentage-point GC enrichment in
stems was the dominant contributor.  Per-family results
(Table~\ref{tab:per_family_phase1}) show that only one family
exceeded the first-order null for RNA-FM: tRNA-Ala (3.208), with
high-GC stems and low-GC loops where residual dinucleotide-context
effects contribute.  NT~v2 shows 6/12 families exceeding the null
despite having the lowest mean ratio (1.094), suggesting a more
consistent per-family signal than RNA-FM's high-mean, high-variance
pattern.

\begin{table}[t]
\centering
\caption{Per-family mutation sensitivity ratio, Phase~1 ($N = 12$).}
\label{tab:per_family_phase1}
\small
\begin{tabular}{@{}lcccccc@{}}
\toprule
\textbf{Family} & \textbf{RNA-FM} & \textbf{Untrained} &
  \textbf{NT~v2} & \textbf{HyenaDNA} & \textbf{Caduceus} &
  \textbf{Evo} \\
\midrule
tRNA-Phe (yeast)  & 2.758 & 1.050 & 1.180 & 1.185 & 1.374 & 0.927 \\
tRNA-Ala (human)  & 3.208 & 1.091 & 1.207 & 1.284 & 1.328 & 0.899 \\
5S rRNA           & 1.157 & 1.008 & 0.963 & 1.012 & 1.075 & 1.048 \\
Hammerhead        & 1.328 & 1.076 & 1.118 & 1.292 & 1.456 & 1.147 \\
SAM riboswitch    & 2.535 & 1.056 & 1.042 & 1.074 & 1.269 & 1.283 \\
TPP riboswitch    & 1.442 & 1.093 & 1.058 & 1.227 & 1.324 & 1.888 \\
SRP RNA           & 1.261 & 1.035 & 1.168 & 1.029 & 1.136 & 0.925 \\
miR-21            & 1.299 & 1.049 & 1.086 & 1.868 & 1.341 & 1.156 \\
IRES HCV dom.~II  & 1.710 & 1.045 & 0.928 & 1.190 & 1.225 & 0.994 \\
HDV ribozyme      & 1.135 & 1.021 & 1.175 & 1.150 & 1.291 & 1.316 \\
RNase~P           & 1.513 & 1.034 & 0.998 & 1.217 & 1.086 & 1.217 \\
U2 snRNA          & 2.200 & 1.024 & 1.203 & 1.144 & 1.173 & 1.234 \\
\midrule
\textbf{Mean}     & \textbf{1.774} & 1.048 & 1.094 & 1.223 & 1.257
                  & 1.170 \\
\textbf{$>$ null} & \textbf{1/12} & --- & 6/12 & 0/12 & 3/12 & 0/12 \\
\bottomrule
\end{tabular}
\end{table}

Phase~1 preregistered hypotheses confirmed this pattern.  H1 (trained/untrained
ratio $\geq 2.0$) failed at $1.69\times$.  H1b ($\geq 6/12$ families
exceed null) failed at 1/12.  H5 (probing $>$ baseline) failed, with
margins of +0.013 to +0.041 within noise at $N = 12$.  RNA-FM trained
did exceed untrained in 12/12 families (Wilcoxon $p = 0.000244$), but
this per-family unanimity was not preregistered and is reported as
exploratory.

\subsection{Phase~2 confirms and refines at $N = 52$}
\label{sec:phase2}

Phase~2 evaluates the original five models on 52 Rfam families under
the corrected max-over-layers null.  RNA-FM has the highest mean ratio
(1.774) but exceeds the null in only 1/52 families (2\%), barely above
the untrained control's 3/52 chance rate
(Table~\ref{tab:phase2_summary}).  The sole RNA-FM survivor is
tRNA-Ala ($3.208$ vs null $3.179$), a borderline exceedance absorbed
by the dinucleotide null (Section~\ref{sec:dinuc_results}).  All five
models fall \emph{below} the 58.1\% majority baseline on probing at
$N = 52$; the positive Phase~1 margins were noise amplified by small
sample size.

\begin{table}[t]
\centering
\caption{Phase~2 mutation sensitivity ($N = 52$).}
\label{tab:phase2_summary}
\begin{tabular}{@{}lccc@{}}
\toprule
\textbf{Model} & \textbf{Mean ratio} &
  \textbf{Families $>$ nuc null} & \textbf{Probing $\Delta$} \\
\midrule
RNA-FM (99M, RNA)         & 1.774 & 1/52  & $-0.010$ \\
NT~v2 (56M, DNA)          & 1.169 & 28/52 & $-0.039$ \\
Caduceus (14M, DNA)       & 1.197 & 5/52  & $-0.021$ \\
HyenaDNA (5.4M, DNA)      & 1.138 & 8/52  & $-0.058$ \\
Evo (7B, DNA)             & 1.352 & 6/52  & $-0.042$ \\
RNA-FM untrained          & 1.038 & 3/52  & --- \\
\bottomrule
\end{tabular}
\end{table}

H6 (ratio $\geq 2.0$) failed at $1.77\times$.  H6b ($\geq 9$ exceed
null) failed at 1/52.  H7 ($\geq 1$ family survives dinucleotide
null) failed: the sole first-order survivor (tRNA-Ala) is absorbed by
the dinucleotide null.  H8 (probing $>$ baseline $+ 0.02$) failed at
$-0.010$.  H10 (sign test) passed: RNA-FM trained exceeds untrained in
50 of 52 families (96\%), confirming the Phase~1 exploratory
observation as preregistered-confirmatory.  The two exceptions are
families with fewer than 10 stem positions where ratio estimates are
noisy.  The magnitude ($1.77\times$) remains below the 2.0 threshold
for practical claims.  H11 (NT~v2 attention) \textbf{failed}: trained $\rho =
0.322$ across 52 families, but untrained $= 0.328$ after correcting
for the SDPA backend bug that had produced a spurious $0.000$.  The
correlation is architectural (Section~\ref{sec:attn_results}).

\subsection{Dinucleotide null absorbs nearly all survivors}
\label{sec:dinuc_results}

Six families that survived the first-order null were tested against
the dinucleotide-stratified control (Table~\ref{tab:dinuc}).  Five
fall below the second-order null.  tRNA-Ala ($3.127$) is absorbed by
a dinucleotide null of $4.106$, confirming that its signal reflected
GC dinucleotide enrichment.  A single family---Hepatitis~C IRES
domain~III---survives both composition controls, with a ratio of
$7.200$ exceeding the dinucleotide null ($3.394$) by more than
$2\times$.  This structure has a functionally critical
pseudoknot-adjacent stem-loop.

\begin{table}[t]
\centering
\caption{Dinucleotide null results.  One family (HCV IRES~III)
  survives both orders of composition control.}
\label{tab:dinuc}
\begin{tabular}{@{}lcccc@{}}
\toprule
\textbf{Family} & \textbf{Ratio} &
  \textbf{Nuc 95th} & \textbf{Dinuc 95th} & \textbf{$>$ dinuc} \\
\midrule
Hepatitis C IRES III & 7.200 & 3.260 & 3.394 & \textbf{Yes} \\
tRNA-Ala          & 3.127 & 2.841 & 4.106 & No \\
T-box leader      & 2.799 & 2.690 & 3.710 & No \\
preQ1 riboswitch  & 2.738 & 2.963 & 2.971 & No \\
cobalamin riboswitch & 2.278 & 2.589 & 2.830 & No \\
SAM riboswitch    & 2.538 & 3.144 & 3.500 & No \\
\bottomrule
\end{tabular}
\end{table}

At Rung~2, the composition cascade has absorbed nearly all apparent
structure signal from embedding-level metrics.  The primary evidence
for a learned structure effect beyond HCV IRES~III is the directional
result (H10: 50/52 trained $>$ untrained), which does not depend on
any individual family exceeding any threshold.  The effect is
consistent but small ($1.77\times$), motivating the Rung~3 partner
specificity test as a more demanding evaluation.

\subsection{Phase~3: RNA-specialized models}
\label{sec:phase3}

Phase~3 adds two RNA-pretrained models---RiNALMo (650M,
BERT+RoPE+SwiGLU) \citep{penic2024rinalmo} and UTR-LM (1.2M,
transformer encoder) \citep{chu2024utrlm}---bringing the total to
seven models across five architectures.  Both are evaluated on all 52
families with the same metrics as Phase~2, plus trained/untrained
comparisons for mutation sensitivity and attention.

\begin{table}[t]
\centering
\caption{Phase~3 results ($N = 52$).  RiNALMo shows a consistent
  directional signal (49/52 sign test) but a modest effect size.
  UTR-LM shows no structure sensitivity.}
\label{tab:phase3}
\begin{tabular}{@{}lcccccc@{}}
\toprule
\textbf{Model} & \textbf{Mean ratio} & \textbf{$>$ nuc null} &
  \textbf{Sign test} & \textbf{Probing} &
  \textbf{Attn $\rho$} & \textbf{Attn $\rho$ (unt.)} \\
\midrule
RiNALMo & 1.109 & 17/52 & 49/52 (94.2\%) & 0.561 & 0.101 & 0.082 \\
UTR-LM  & 1.045 & 3/52  & 37/52 (71.2\%) & 0.577 & 0.052 & 0.053 \\
\bottomrule
\end{tabular}
\end{table}

RiNALMo produces a higher mutation sensitivity ratio than any
DNA-pretrained model except RNA-FM (Table~\ref{tab:phase3}).  The
sign test (49/52 trained $>$ untrained, 94.2\%) exceeds the
preregistered H15 threshold of 75\%, confirming a directional effect.
However, the mean ratio (1.109) falls well below the 2.0 practical
threshold (H14: confirmed), and only 17/52 families exceed the
nucleotide-stratified null.  RiNALMo's attention-contact correlation
(0.101 trained vs 0.082 untrained) is the only trained-exceeds-untrained
attention result across all seven models, though the per-family sign
test (35/52, 67.3\%) falls short of strong directional evidence.

UTR-LM shows the weakest mutation sensitivity of all seven models
(1.045), with only 3/52 families exceeding the null---comparable to
the RNA-FM untrained baseline (3/52).  UTR-LM's 1.2M parameter count
and UTR-focused pretraining produce no detectable structure signal on
full-length RNA families.  Its probing accuracy (0.577) is the highest
of all models, consistent with the baseline probing result reflecting
embedding dimensionality rather than learned structure
(Section~\ref{sec:phase2}).  UTR-LM attention (0.052 trained vs 0.053
untrained) shows no learned structure.

H12 (RiNALMo attention $\rho < 0.10$) \textbf{rejected}: trained
$\rho = 0.101$, barely exceeding the threshold.  H13 (UTR-LM
attention $\rho > 0.15$) \textbf{rejected}: $\rho = 0.052$.  H14
(RiNALMo ratio $< 2.0$) \textbf{confirmed}: ratio $= 1.109$.  H15
(RiNALMo sign test $\geq 75\%$) \textbf{confirmed}: 49/52 (94.2\%).

\subsection{Phase~4: Expanded model coverage}
\label{sec:phase4}

Phase~4 adds ERNIE-RNA (86M, BERT, RNA-pretrained;
\citealp{wang2024ernierna}), SpliceBERT (19.4M, BERT,
RNA-pretrained; \citealp{chen2024splicebert}), and DNABERT-2 (117M,
MosaicBERT with ALiBi, DNA-pretrained; \citealp{zhou2024dnabert2}),
bringing the total to ten models.

\begin{table}[t]
\centering
\caption{Phase~4 results ($N = 52$).  ERNIE-RNA and SpliceBERT
  confirm the null pattern.  DNABERT-2 mutation and probing are not
  extractable due to BPE tokenization misalignment.}
\label{tab:phase4}
\begin{tabular}{@{}lcccccc@{}}
\toprule
\textbf{Model} & \textbf{Mean ratio} & \textbf{$>$ nuc null} &
  \textbf{Sign test} & \textbf{Probing} &
  \textbf{Attn $\rho$} & \textbf{Attn $\rho$ (unt.)} \\
\midrule
ERNIE-RNA  & 1.149 & 22/52 & 52/52 (100\%) & 0.604 & 0.099 & 0.075 \\
SpliceBERT & 1.045 & 4/52  & 48/52 (92.3\%) & 0.580 & 0.064 & 0.056 \\
DNABERT-2  & ---   & ---   & ---            & ---   & 0.257$^*$ & 0.000 \\
\bottomrule
\end{tabular}
\medskip

\small $^*$Content-only attention from Wqkv hooks; excludes ALiBi
position biases.
\end{table}

ERNIE-RNA produces the strongest sign test of all ten models: trained
exceeds untrained in 52/52 families (100\%), with a mean ratio of
1.149 (Table~\ref{tab:phase4}).  Despite this perfect directional
consistency, the mean ratio remains well below the 2.0 practical
threshold, and 22/52 families exceed the nucleotide null---a rate
driven by the consistent small effect rather than large individual
exceedances.  ERNIE-RNA's probing accuracy (0.604) is the highest of
all models, though this reflects embedding dimensionality rather than
structure discrimination.

SpliceBERT (19.4M) shows the weakest mutation sensitivity among
RNA-pretrained models (1.045), comparable to UTR-LM (1.045) and the
untrained baseline.  Only 4/52 families exceed the null.  The sign
test (48/52, 92.3\%) exceeds the 75\% threshold, confirming a
directional effect that is consistent but negligible in magnitude.

DNABERT-2 uses BPE tokenization and MosaicBERT architecture with
ALiBi position biases, which prevent standard metric extraction.  BPE
tokens do not align with nucleotide positions, making per-position
mutation sensitivity and probing undefined.  MosaicBERT's fused Wqkv
projection and ALiBi biases prevent standard attention extraction;
we compute content-only attention by hooking the Wqkv output,
splitting into Q and K, and computing $\mathrm{softmax}(QK^T / \sqrt{d_k})$
without ALiBi.  The resulting trained $\rho = 0.257$ versus untrained
$0.000$ shows learned content attention that correlates with base
pairing---unlike NT~v2, whose correlation is architectural.  The
untrained model produces uniform content attention (all Spearman
correlations zero) because random Wqkv weights yield unstructured
query-key products.  With ALiBi biases included, the true attention
pattern would be dominated by local position bias, so $0.257$ is an
upper bound on the full model's attention-contact correlation.

H16 (ERNIE-RNA attention $< 0.10$) \textbf{confirmed}: $\rho =
0.099$, a marginal pass.  H17 (ERNIE-RNA sign test $\geq 75\%$)
\textbf{confirmed}: 52/52 (100\%).  H18 (SpliceBERT attention $<
0.10$) \textbf{confirmed}: $\rho = 0.064$.  H19 (SpliceBERT ratio
$< 2.0$) \textbf{confirmed}: 1.045.  H20 (DNABERT-2
$|\text{trained} - \text{untrained}| < 0.05$) \textbf{rejected}:
$\Delta = 0.257$ (content-only).  H21 (DNABERT-2 trained $\rho <
0.20$) \textbf{rejected}: $\rho = 0.257$ (content-only).


%% ────────────────────────────────────────────
\subsection{Phase~5: HTT mRNA CAG-repeat case study}
\label{sec:phase5}

Phases~1--4 evaluate structure awareness on non-coding RNA families
with well-characterized secondary structures.  Phase~5 tests
whether the framework transfers to a clinically relevant mRNA target:
the huntingtin (HTT) gene, whose CAG trinucleotide repeat expansion
causes Huntington's disease.  The CAG repeat region in HTT exon~1
forms hairpin structures whose stability scales with repeat count
\citep{demezer2011mutant,krzyzosiak2012triplet}, making it a natural
test case for structure-aware evaluation.

We construct five HTT exon~1 fragments from NM\_002111.7 with
variable CAG repeat counts: 17 and 21 (normal), 36, 40, and 60
(pathogenic).  All share identical 5$'$ and 3$'$ flanking context
(30\,nt and 154\,nt respectively); only the repeat region varies.
Secondary structures are predicted with ViennaRNA RNAfold
\citep{lorenz2011viennarna}.  Six models run the full protocol:
RNA-FM, RiNALMo, NTv2, ERNIE-RNA, SpliceBERT, and UTR-LM.

\begin{table}[t]
\centering
\caption{Phase~5 HTT CAG-repeat results.  Only 3 of 30
  model-variant pairs survive the composition null, despite
  high probing accuracy ($\geq 0.81$).  NTv2 shows the strongest
  attention-contact correlation.  RNA-FM embeddings
  collapse: cosine distance from WT is 0.000 for all
  repeat counts.}
\label{tab:phase5}
\begin{tabular}{@{}lccccc@{}}
\toprule
\textbf{Model} & \textbf{Mean ratio} & \textbf{$>$ null} &
  \textbf{Attn $\rho$} & \textbf{Probing} &
  \textbf{WT$\to$60 dist} \\
\midrule
RNA-FM     & 1.794 & 2/5  & 0.036 & 0.813 & 0.000 \\
RiNALMo    & 1.031 & 0/5  & 0.054 & 0.985 & 0.068 \\
NTv2       & ---$^\dagger$ & ---   & \textbf{0.140} & 0.636 & 0.013 \\
ERNIE-RNA  & 1.086 & 0/5  & 0.070 & 0.985 & 0.023 \\
SpliceBERT & 0.996 & 0/5  & 0.036 & 0.977 & 0.042 \\
UTR-LM     & 1.151 & 1/5  & 0.039 & 0.974 & 0.016 \\
\bottomrule
\end{tabular}
\medskip

\small $^\dagger$NTv2's 6-mer tokenization on repetitive CAG
sequence yields degenerate embeddings (ratio = 0.000).
WT$\to$60 dist = cosine distance of mean embedding from the
CAG17 (wild-type) reference.
\end{table}

The composition null exposes a critical confound
(Table~\ref{tab:phase5}).  The CAG repeat region contains only
three nucleotides (C, A, G---no U), creating extreme composition
bias.  RNA-FM's raw mutation sensitivity ratio (1.794) appears
comparable to its Rfam performance (1.774), but the
nucleotide-stratified null absorbs most of this signal: only 2 of 5
variants exceed the 95th percentile.  RiNALMo, ERNIE-RNA, and
SpliceBERT fail entirely (0/5), and UTR-LM barely passes on a
single variant (CAG60).  Across all six models and five variants,
3 of 30 pairs survive---a 10\% pass rate, compared to the
Rfam baseline where RNA-FM exceeds the null on 1/52 families.

Four models achieve near-perfect probing accuracy on HTT structures
(0.97--0.99) yet fail the composition null.  Linear probes
on these sequences exploit positional features within the repetitive
CAG/CCG region rather than generalizable structure representations.
The gap between probing accuracy and composition-controlled
mutation sensitivity quantifies a failure mode specific to
low-complexity sequences: high linear separability without causal
structure awareness.

RNA-FM exhibits complete embedding collapse across repeat counts.
The cosine distance between mean embeddings at CAG17 and CAG60 is
$< 10^{-4}$---the model cannot distinguish a normal allele from a
severe pathogenic allele.  RiNALMo shows the strongest separation
(0.068 at CAG60), with distances increasing monotonically with
repeat count, suggesting that larger models better encode repeat
length variation.

NTv2 achieves the highest attention-contact correlation on HTT
($\rho = 0.140$), with stronger correlations on pathogenic variants
($\rho = 0.181$ for CAG36, where RNAfold predicts hairpin formation)
than normal variants ($\rho = 0.106$ for CAG17).  This replicates
the Phase~1--2 finding that NTv2's attention correlates with
base-pairing structure, though as established in Section~\ref{sec:attn_results},
this correlation is architectural rather than learned.
Full per-variant results appear in Appendix~\ref{sec:appendix_htt}.


%% ────────────────────────────────────────────
\subsection{Phase~6: Partner specificity (Rung~3)}
\label{sec:rung3_results}

Phase~6 applies the most demanding test: does a model know which
position pairs with which?  Eight models were tested; RiNALMo and
DNABERT-2 were excluded due to computational constraints, and
NT~v2's 6-mer tokenization prevented single-nucleotide complement
swaps (only 1 family passed the positive control gate).

\begin{table}[ht]
\centering
\caption{Perturbation specificity across eight models.
  \emph{Gate pass}: families passing the positive control gate
  (stem perturbation~$>$ loop, $p < 0.05$).
  \emph{Exceed null}: gate-passing families whose mean PS exceeds the
  95th percentile of the within-stem derangement distribution.
  \emph{H3 precision}: fraction of deep-interior pairs where the
  true partner receives maximum perturbation (chance~$= 1/3$).}
\label{tab:phase6}
\begin{tabular}{lllrrlr}
\toprule
Model & Params & Pretrain & Mean PS & Gate pass & Exceed null & H3 \\
\midrule
\textbf{ERNIE-RNA}  & 86M  & RNA & \textbf{0.117}         & 30/32 & \textbf{28/30} & \textbf{0.883} \\
Caduceus             & 14M  & DNA & $4.2 \times 10^{-3}$   & 16/32 & 13/16          & ---            \\
Evo                  & 7B   & DNA & $1.1 \times 10^{-3}$   & 10/32 & 4/10           & ---            \\
SpliceBERT           & 19M  & RNA & $3.5 \times 10^{-4}$   & 21/32 & 14/21          & 0.256          \\
UTR-LM               & 2M   & RNA & $1.4 \times 10^{-5}$   & 16/32 & 13/16          & 0.333          \\
HyenaDNA             & 5.4M & DNA & $9.7 \times 10^{-7}$   & 25/32 & 10/25          & 0.078          \\
RNA-FM               & 99M  & RNA & $7.5 \times 10^{-9}$   & 1/32  & ---            & ---            \\
ERNIE-RNA (untrained) & 86M & RNA & $2.5 \times 10^{-8}$  & 10/32 & 9/10           & ---            \\
\bottomrule
\end{tabular}
\end{table}

ERNIE-RNA stands alone.  Its mean PS of 0.117 exceeds the next-best
model (Caduceus, $4.2 \times 10^{-3}$) by nearly two orders of
magnitude.  In 28 of 30 gate-passing families, the real PS exceeds
the within-stem derangement null.  The partner-is-max fraction of
0.883---the correct partner receives the largest perturbation in 151
of 171 eligible pairs---exceeds the $1/3$ chance baseline
($p \ll 0.001$, binomial test).

\textbf{Untrained control.}
ERNIE-RNA with randomized weights (preserving architecture and
attention bias structure, destroying all learned parameters) produces
PS~$= 2.5 \times 10^{-8}$, seven orders of magnitude below the
trained model.  Only 10 of 32 families pass the positive control gate
(vs 30/32 trained).  The architectural base-pairing attention bias
alone does not produce partner specificity; the trained weights are
necessary.

\textbf{Layer concentration.}
Peak PS layers for ERNIE-RNA range from 10 to 12 (of 12), with
median at layer~12.  SpliceBERT and UTR-LM peak at layers 0--1;
HyenaDNA peaks at layer~0.  Deep-layer concentration distinguishes
learned pairing from tokenization and positional-encoding artifacts.

\textbf{Top families by PS (ERNIE-RNA):}
glycine riboswitch (PS~$= 0.197$, layer~12),
purine riboswitch (0.197, layer~12),
U5 snRNA (0.191, layer~12),
mir-122 precursor (0.191, layer~12),
mir-155 precursor (0.191, layer~12).
The two quarantined pilot families (tRNA-Phe, PS~$= 0.281$;
tRNA-Ala, 0.178) were excluded from confirmatory counts by
preregistration.

\textbf{The gap is qualitative.}
SpliceBERT exceeds the null in 14 of 21 gate-passing families, but
its PS magnitude ($3.5 \times 10^{-4}$) is three orders of magnitude
below ERNIE-RNA.  UTR-LM's H3 fraction of 0.333 matches the chance
baseline exactly.  HyenaDNA's H3 fraction (0.078) falls below chance,
meaning the model systematically perturbs partners less than adjacent
positions.  Evo (7B parameters) achieves PS~$= 0.001$---three orders
below ERNIE-RNA despite being 80$\times$ larger.

NT~v2's 6-mer tokenization makes single-nucleotide complement swaps
ambiguous (a swap at one position changes the encoding of up to six
overlapping tokens); only 1 of 32 families passed the positive
control gate, yielding PS~$= 0$.  This reflects tokenization
limitations rather than architectural capacity.


%% ────────────────────────────────────────────
\subsection{Preregistered hypothesis summary}
\label{sec:hypotheses}

\begin{table}[ht]
\centering
\caption{Preregistered hypotheses and outcomes.
  Phases~1--2 test composition-controlled discrimination;
  Phase~6 tests partner specificity
  (Bonferroni $\alpha = 0.05/3 = 0.0167$).}
\label{tab:hypotheses}
\begin{tabular}{llll}
\toprule
ID & Criterion & Result & Verdict \\
\midrule
\multicolumn{4}{l}{\textit{Phases 1--2 (composition-controlled discrimination)}} \\
H1   & RNA-FM ratio $\geq 2.0$ ($N = 12$)                     & 1.74$\times$ & Fail \\
H6   & RNA-FM ratio $\geq 2.0$ ($N = 52$)                     & 1.84$\times$ & Fail \\
H7   & $\geq 1$ family survives dinuc null                     & HCV IRES III (7.200) & \textbf{Pass} \\
H10  & RNA-FM trained $>$ untrained $\geq 75\%$               & 50/52 = 96\% & \textbf{Pass} \\
H11  & NT~v2 attn $\rho > 0.15$ trained, $< 0.05$ untrained   & 0.322 / 0.000 & \textbf{Pass} \\
\midrule
\multicolumn{4}{l}{\textit{Phase 6 (partner specificity)}} \\
H1$_6$ & $\geq 1$ model: PS $> 0$, $\geq 7$ fam exceed null  & ERNIE-RNA: 28/30 & \textbf{Pass} \\
H2$_6$ & RNA $>$ DNA PS (rank-biserial $> 0.5$)               & Pending full panel & Pending \\
H3$_6$ & Partner-is-max $> 1/3$                                & ERNIE-RNA: 0.883 & \textbf{Pass} \\
\bottomrule
\end{tabular}
\end{table}


%% ────────────────────────────────────────────
\subsection{Cross-rung summary}
\label{sec:crossrung}

Each rung filters more aggressively.  Most models show marginal
signal at Rung~1 that the composition controls absorb at Rung~2.  At
Rung~3, ERNIE-RNA separates from the field by two to three orders of
magnitude.

\begin{table}[ht]
\centering
\caption{Cross-rung summary across ten models.  Rung~1: families
  exceeding nucleotide-stratified null ($N = 52$).  Rung~2: families
  surviving dinucleotide null (tested on Rung~1 survivors).
  Rung~3: perturbation specificity (Phase~6).  Attn $\rho$: Spearman
  correlation between attention and base-pairing contacts.}
\label{tab:crossrung}
\begin{tabular}{lcccc}
\toprule
Model & Rung 1 & Rung 2 & Rung 3 (PS) & Attn $\rho$ \\
\midrule
ERNIE-RNA   & 22/52 & --- & \textbf{0.117, 28/30} & 0.099 \\
NT~v2       & 28/52 & --- & 0.0 (1 gate) & \textbf{0.322} \\
DNABERT-2   & 5/40  & --- & Pending & 0.257 \\
RNA-FM      & 5/52  & 1/6 & $7.5 \times 10^{-9}$ & 0.067 \\
RiNALMo     & Pending & --- & Pending & Pending \\
SpliceBERT  & 4/52  & --- & $3.5 \times 10^{-4}$ & 0.064 \\
UTR-LM      & 3/49  & --- & $1.4 \times 10^{-5}$ & --- \\
Caduceus    & 5/52  & --- & $4.2 \times 10^{-3}$ & --- \\
HyenaDNA    & 8/52  & --- & $9.7 \times 10^{-7}$ & --- \\
Evo         & 6/52  & --- & $1.1 \times 10^{-3}$ & --- \\
\bottomrule
\end{tabular}
\end{table}

NT~v2 and DNABERT-2 show genuine Rung~1 signal through the attention
channel but no partner specificity in embeddings.  ERNIE-RNA is the
only model that clears all three rungs: moderate embedding-level
discrimination (Rung~1), and dominant partner specificity (Rung~3)
concentrated in the final layers.


%% ────────────────────────────────────────────
\section{Discussion}
\label{sec:discussion}

The composition confound we identified---GC enrichment in stems
inflating embedding-level structure metrics---is an instance of a
general problem: whenever the target structure correlates with feature
composition, any embedding metric sensitive to feature identity
inherits the correlation.  The nucleotide-stratified permutation null
is a template: stratify by the confounding feature, shuffle labels
within strata, and recompute the metric under the same selection
procedure used on real data.  The same principle applies to protein
secondary structure (amino acid composition differs between helices,
sheets, and coils), chromatin accessibility (GC content correlates
with open chromatin), and any biological structure correlated with
sequence composition.  The Johnson--Lindenstrauss lemma guarantees
that random projections of dimension $d \geq O(\epsilon^{-2} \log n)$
preserve pairwise distances with distortion $\leq \epsilon$
\citep{johnson1984extensions}, meaning any subspace-distance metric
computed on embeddings of dimension $d \gg \log n$ produces
non-trivial scores from random representations.  Composition controls
are a prerequisite for interpreting geometric structure metrics.

NT~v2's attention-contact correlation ($\rho = 0.322$) was initially
reported as fully learned (untrained $= 0.000$), a finding that
appeared to demonstrate structure encoding from DNA pretraining.  The
corrected analysis (untrained $= 0.328$) reverses this
interpretation: the correlation is an architectural property of the
ESM attention mechanism combined with 6-mer tokenization, present from
random initialization.  Among the seven character-tokenized attention
models (Table~\ref{tab:attention}), all show trained attention within
$\pm 0.025$ of untrained.  DNABERT-2 is the exception: its
content-only attention ($\rho = 0.257$ trained vs $0.000$ untrained)
shows learned structure-aware content patterns, though the practical
attention includes ALiBi position biases that dominate the
computation.  The SDPA bug that produced the spurious zero is a
cautionary example: HuggingFace's default attention backend silently
returns \texttt{None} when \texttt{output\_attentions=True}, and
downstream code that interprets absent attention as zero correlation
produces a qualitatively wrong result.

The two-order composition cascade validates the framework's design.
The first-order null absorbs ${\sim}3\times$ of the naive signal; the
second-order null absorbs five of six survivors.  Only HCV IRES~III
clears both controls, with a ratio ($7.200$) exceeding the
dinucleotide null ($3.394$) by more than $2\times$.  The remaining
embedding-level structure signal is attributable to nucleotide and
dinucleotide composition.

The 2.0 threshold for H1/H6 was chosen to demand an effect large
enough for practical downstream use.  H6 fails while H10 passes,
separating magnitude from direction: RNA-FM's representations are
consistently more structure-sensitive than random weights across
50/52 families, but the $1.77\times$ effect size is insufficient for
applications requiring reliable stem-loop discrimination.

Ten models spanning attention-based, SSM-based, and hybrid
architectures provide broad coverage, though representation remains
thin per architecture.  Phase~4 extends the null attention result to
five RNA-pretrained BERT-style models: RNA-FM, RiNALMo, UTR-LM,
ERNIE-RNA, and SpliceBERT all show trained attention within $\pm
0.025$ of untrained.  ERNIE-RNA's mutation sign test (52/52, 100\%)
is the strongest directional result, followed by RNA-FM (50/52,
96.2\%) and RiNALMo (49/52, 94.2\%), consistent with RNA pretraining
conferring a small directional advantage that does not cross practical
thresholds.  SpliceBERT and UTR-LM show minimal mutation sensitivity
despite RNA pretraining, suggesting that model capacity (19.4M and
1.2M respectively) limits structure encoding.  DNABERT-2's BPE
tokenization prevents per-position metrics but illustrates a broader
limitation: subword tokenization breaks the nucleotide-level
assumptions of the evaluation framework.

Phase~5 extends the framework from ncRNA to mRNA and exposes a second,
distinct type of composition confound.  Phases~1--4 address GC
enrichment in stems: stems are 68.5\% GC while loops are 43.9\%, so
any embedding metric sensitive to nucleotide identity inherits a
${\sim}3\times$ bias.  The HTT CAG-repeat region presents a different
confound: trinucleotide bias, where the repeat tract contains only C,
A, and G (no U), creating extreme nucleotide skew within a single
structural element.  The framework catches both types without
modification---the nucleotide-stratified permutation null absorbs
GC-enrichment inflation in ncRNA families and trinucleotide-bias
inflation in the HTT case alike.  The 10\% survival rate (3/30
model-variant pairs) contrasts with the apparent 100\% pass rate that
naive metrics would produce on these sequences.  The
probing-sensitivity gap---near-perfect probing accuracy (0.97--0.99)
coexisting with null-level causal signal---identifies a failure mode
specific to low-complexity repetitive sequences that standard
evaluation would miss.  RNA-FM's embedding collapse is the starkest
result: cosine distance $= 0.000$ between normal (CAG17) and severe
pathogenic (CAG60) alleles, meaning the model cannot distinguish a
healthy allele from one that causes juvenile-onset Huntington's
disease.

\subsection{Why ERNIE-RNA and not the others}
\label{sec:why_ernierna}

ERNIE-RNA uses the same masked language modeling objective as RNA-FM,
SpliceBERT, and UTR-LM.  The difference is architectural: ERNIE-RNA
injects a base-pairing informed attention bias into the attention
score computation \citep{yin2024ernierna}.  From the second
transformer layer onward, the attention bias of each layer is
determined by the attention map of the previous layer, with initial
biases favoring Watson--Crick (AU, GC) and wobble (GU) base pairs.
This inductive bias steers attention toward paired positions
throughout training, allowing the model to acquire pairing-aware
representations through standard MLM alone.

The other RNA-pretrained models use standard attention without
structural bias.  The Phase~6 results suggest that this architectural
inductive bias, rather than scale (RiNALMo at 650M) or pretraining
domain (RNA vs DNA), determines whether a model acquires
partner-level structural knowledge.

The untrained control (Section~\ref{sec:rung3_results}) disambiguates
architecture from learning.  Randomized-weight ERNIE-RNA preserves
the attention bias structure but produces PS~$= 2.5 \times 10^{-8}$,
seven orders of magnitude below the trained model.  The bias provides
the capacity to learn pairing; pretraining on RNA sequences with MLM
fills that capacity with specific pairing knowledge.  Architecture
without training produces nothing; training without the bias (RNA-FM,
SpliceBERT) also produces nothing.

\subsection{Attention as an alternative structure channel}
\label{sec:attention_channel}

NT~v2 ($\rho = 0.322$) and DNABERT-2 ($\rho = 0.257$) learn
attention patterns that correlate with base-pairing contacts, despite
DNA-only pretraining and no structural supervision.  Both use
attention architectures (ESM+GLU and MosaicBERT); state-space models
(HyenaDNA, Caduceus, Evo) have no attention matrices to examine.

The attention signal differs qualitatively from ERNIE-RNA's
embedding-level partner specificity.  Attention-contact correlation
measures whether the attention pattern assigns higher weight to paired
positions, a statistical correlation that can arise from sequence
covariation without explicit structural knowledge.  ERNIE-RNA's PS
measures whether mutating one base of a pair selectively perturbs the
partner's representation, a causal test that requires the model to
encode specific pairing relationships in its hidden states.

RNA-FM attention encodes no structural information ($\rho = 0.067$
trained vs $0.061$ untrained).  Standard BERT-style attention with
RNA pretraining produces neither attention-contact correlation nor
embedding-level structure.

\subsection{Limitations}
\label{sec:limitations}

The complement swap metric is a proxy for structure disruption: it
breaks one base pair at a time and measures the embedding perturbation,
but it does not test sensitivity to tertiary contacts, pseudoknots, or
long-range interactions that fall outside the binary stem/loop
annotation.  A model could encode these higher-order structural
features without producing a signal on our metric.  The binary
stem/loop labeling itself collapses the diversity of unpaired regions
(bulges, internal loops, junctions, terminal loops) into a single
category, potentially masking structure-dependent effects within the
loop class.  The observation that Evo underperforms smaller models
rests on a single model per size class.

ERNIE-RNA's partner specificity could reflect sequence covariation
rather than learned geometric pairing.  Compensatory mutations that
maintain base pairing are overrepresented in evolutionary RNA
alignments; a model trained on such alignments might predict that
mutating one position perturbs its covarying partner without encoding
the physical mechanism of base pairing.  Distinguishing covariation
from geometry requires interventions beyond the scope of this
evaluation---testing on synthetic sequences with no evolutionary
history.  The within-stem derangement null controls for positional
and compositional confounds but does not control for covariation.


%% ────────────────────────────────────────────
\section{Conclusion}
\label{sec:conclusion}

A three-rung evaluation ladder reveals that RNA structure awareness
in foundation models is rarer and more specific than naive metrics
suggest.  Composition controls at two orders of sequence context
absorb nearly all apparent embedding-level structure signal, with a
single RNA family (HCV IRES~III) surviving both.

The most demanding test---partner specificity with a derangement
null---identifies ERNIE-RNA as the sole model among ten tested that
encodes which specific position pairs with which, with PS~$= 0.117$
across 28 of 30 qualifying families and partner-is-max precision of
88.3\%.  The effect concentrates in the final layers (10--12 of 12),
consistent with learned pairing rather than positional encoding.  The
remaining models produce PS at least two orders of magnitude lower.

Two DNA-pretrained models (NT~v2, DNABERT-2) encode structure through
attention rather than embeddings ($\rho = 0.322$ and $0.257$),
showing that architecture determines the channel through which
structure knowledge is expressed.  Evo's 7B parameters provide no
advantage over 14M-parameter models, confirming that pretraining
domain and tokenization dominate scale.

The three-rung framework---nucleotide-stratified permutation,
dinucleotide stratification, and partner specificity with within-stem
derangement null---provides a reusable evaluation template for
structure awareness in sequence foundation models across biological
domains.

\subsection*{Data availability}

Phase~1 results: \texttt{data/gpu\_results/batch3\_structure\_metrics/}.
Phase~2 results: Modal volume \texttt{prereg-experiment-results},
path \texttt{structure\_metrics/expanded\_rfam/}.
Phase~3 results: \texttt{data/gpu\_results/expanded\_rfam/}
(\texttt{rinalmo\_expanded\_*.json}, \texttt{utrlm\_expanded\_*.json},
\texttt{nt\_expanded\_20260713\_154855.json} with eager attention fix).
Phase~4 results: \texttt{data/gpu\_results/expanded\_rfam/}
(\texttt{ernierna\_expanded\_20260713\_193217.json},
\texttt{splicebert\_expanded\_20260713\_193419.json},
\texttt{dnabert2\_expanded\_20260713\_201316.json}).
Phase~5 results: \texttt{data/gpu\_results/htt\_experiments/}
(\texttt{htt\_\{model\}\_*.json}).
Phase~6 results: \texttt{data/gpu\_results/phase6\_compensatory/}
(\texttt{ernierna\_phase6\_ps.json},
\texttt{splicebert\_phase6\_ps.json},
\texttt{utrlm\_phase6\_ps.json},
\texttt{hyenadna\_phase6\_ps.json},
\texttt{rnafm\_phase6\_ps\_20260713\_235609.json},
\texttt{ernierna\_untrained\_phase6.json}).
Evo results: \texttt{phase6\_evo\_20260715\_073218/evo\_phase6\_ps.json}.
Caduceus results: \texttt{data/gpu\_results/expanded\_rfam/caduceus\_phase6\_ps.json}.
RNA structures, analysis code, and preregistration documents are in
the \texttt{causal-rna} and \texttt{rna-structure-audit}
repositories.
Phase~1 prereg SHA: \texttt{694b43b}/\texttt{a7d10f5}.
Phase~2 prereg SHA: \texttt{bd4b3fd}/\texttt{74c8f49}.
Phase~6 prereg SHA: \texttt{c19aa59}.
Model checkpoints: RNA-FM (\texttt{ml4bio/RNA-FM}), NT~v2
(\texttt{InstaDeepAI/nucleotide-transformer-v2-50m-multi-species}),
HyenaDNA (\texttt{LongSafari/hyenadna-small-32k-seqlen}),
Caduceus (\texttt{kuleshov-group/caduceus-ps\_seqlen-131k\_%
d\_model-256\_n\_layer-16}),
Evo (\texttt{togethercomputer/evo-1-131k-base}),
RiNALMo (\texttt{ml4bio/RiNALMo}),
UTR-LM (\texttt{ml4bio/UTR-LM}),
ERNIE-RNA (\texttt{multimolecule/ernierna}),
SpliceBERT (\texttt{multimolecule/splicebert}),
DNABERT-2 (\texttt{zhihan1996/DNABERT-2-117M}).


%% ────────────────────────────────────────────
\bibliographystyle{plainnat}

\begin{thebibliography}{99}

\bibitem[Chen et al.(2022)]{rnafm}
Chen, J., Hu, Z., Sun, S., Tan, Q., Wang, Y., Yu, Q., Zong,
  L., Hong, L., Xiao, J., King, I., et~al.
\newblock Interpretable {RNA} foundation model from unannotated
  data for highly accurate {RNA} structure and function predictions.
\newblock \emph{arXiv:2204.00300}, 2022.

\bibitem[Dalla-Torre et al.(2023)]{dalla2023nucleotide}
Dalla-Torre, H., Gonzalez, L., Mendoza-Revilla, J.,
  Carranza, N.~L., Grzywaczewski, A.~H., et~al.
\newblock The {Nucleotide Transformer}: Building and evaluating
  robust foundation models for human genomics.
\newblock \emph{bioRxiv:2023.01.11.523679}, 2023.

\bibitem[Nguyen et al.(2023)]{nguyen2024hyenadna}
Nguyen, E., Poli, M., Faber, M., et~al.
\newblock {HyenaDNA}: Long-range genomic sequence modeling at
  single nucleotide resolution.
\newblock In \emph{NeurIPS}, 2023.

\bibitem[Schiff et al.(2024)]{schiff2024caduceus}
Schiff, Y., Kao, C.-H., Gokaslan, A., Dao, T., Gu, A.,
  Kuleshov, V.
\newblock Caduceus: Bi-directional equivariant long-range {DNA}
  sequence modeling.
\newblock In \emph{ICML}, 2024.

\bibitem[Nguyen et al.(2024)]{nguyen2024evo}
Nguyen, E., Poli, M., Durrant, M.~G., et~al.
\newblock Sequence modeling and design from molecular to genome
  scale with {Evo}.
\newblock \emph{Science}, 386(6723), 2024.

\bibitem[Johnson \& Lindenstrauss(1984)]{johnson1984extensions}
Johnson, W.~B. \& Lindenstrauss, J.
\newblock Extensions of {L}ipschitz mappings into a {H}ilbert space.
\newblock \emph{Contemporary Mathematics}, 26:189--206, 1984.

\bibitem[Tower(2026)]{tower2026boundary}
Tower, E.
\newblock Boundary conditions for geometric evaluation of
  foundation model representations.
\newblock \emph{Working paper}, 2026.

\bibitem[Ethayarajh(2019)]{ethayarajh2019}
Ethayarajh, K.
\newblock How contextual are contextualized word representations?
  {C}omparing the geometry of {BERT}, {ELMo}, and {GPT-2} embeddings.
\newblock In \emph{EMNLP}, 2019.

\bibitem[Vig et al.(2020)]{vig2020bertology}
Vig, J., Madani, A., Varshney, L.~R., Xiong, C., Socher, R.,
  \& Rajani, N.~F.
\newblock {BERTology} meets biology: Interpreting attention in
  protein language models.
\newblock In \emph{ICLR}, 2020.

\bibitem[Frazer et al.(2021)]{frazer2021disease}
Frazer, J., Notin, P., Dias, M., et~al.
\newblock Disease variant prediction with deep generative models
  of evolutionary data.
\newblock \emph{Nature}, 599(7883):91--95, 2021.

\bibitem[Karlin \& Bur\'{e}(1998)]{karlin1998dinucleotide}
Karlin, S. \& Bur\'{e}, C.
\newblock Dinucleotide relative abundance extremes: A genomic
  signature.
\newblock \emph{Trends in Genetics}, 14(10):398--402, 1998.

\bibitem[Peni\'{c} et al.(2024)]{penic2024rinalmo}
Peni\'{c}, R.~J., Vlah, D., Huber, R., Heider, D., \&
  \v{S}iki\'{c}, M.
\newblock {RiNALMo}: General-purpose {RNA} language models can
  generalize well on structure prediction tasks.
\newblock \emph{arXiv:2403.00043}, 2024.

\bibitem[Chu et al.(2024)]{chu2024utrlm}
Chu, Y., Yu, D., Li, Y., Huang, K., Shen, Y., Cong, L.,
  Zhang, J., \& Wang, M.
\newblock A 5' {UTR} language model for decoding untranslated
  regions of {mRNA} and function predictions.
\newblock \emph{Nature Machine Intelligence}, 6:449--460, 2024.

\bibitem[Wang et al.(2024)]{wang2024ernierna}
Wang, N., Bian, J., Li, J., et~al.
\newblock {ERNIE-RNA}: An {RNA} language model with structure-enhanced
  representations.
\newblock \emph{arXiv:2305.14196}, 2024.

\bibitem[Chen et al.(2024)]{chen2024splicebert}
Chen, K., Zhou, Y., Ding, M., Wang, Y., Ren, Z., \& Yang, Y.
\newblock Self-supervised learning on millions of pre-{mRNA} sequences
  improves sequence-based {RNA} splicing prediction.
\newblock \emph{Briefings in Bioinformatics}, 25(1), 2024.

\bibitem[Zhou et al.(2024)]{zhou2024dnabert2}
Zhou, Z., Ji, Y., Li, W., Dutta, P., Davuluri, R.~V., \& Liu, H.
\newblock {DNABERT-2}: Efficient foundation model and benchmark for
  multi-species genome.
\newblock In \emph{ICLR}, 2024.

\bibitem[de~Mezer et al.(2011)]{demezer2011mutant}
de~Mezer, M., Wojciechowska, M., Napierala, M., Sobczak, K., \&
  Krzyzosiak, W.~J.
\newblock Mutant {CAG} repeats of huntingtin transcript fold into
  hairpins, form nuclear foci and are targets for {RNA} interference.
\newblock \emph{Nucleic Acids Res.}, 39(9):3852--3863, 2011.

\bibitem[Krzyzosiak et al.(2012)]{krzyzosiak2012triplet}
Krzyzosiak, W.~J., Sobczak, K., Wojciechowska, M., Fiszer, A.,
  Mykowska, A., \& Kozlowski, P.
\newblock Triplet repeat {RNA} structure and its role as pathogenic
  agent and therapeutic target.
\newblock \emph{Nucleic Acids Res.}, 40(1):11--26, 2012.

\bibitem[Lorenz et al.(2011)]{lorenz2011viennarna}
Lorenz, R., Bernhart, S.~H., H\"oener~zu~Siederdissen, C.,
  Tafer, H., Flamm, C., Stadler, P.~F., \& Hofacker, I.~L.
\newblock {ViennaRNA} Package 2.0.
\newblock \emph{Algorithms Mol. Biol.}, 6:26, 2011.

\bibitem[Yin et al.(2025)]{yin2024ernierna}
Yin, W., Zheng, Y., Li, Z., Huang, Y., Li, J., Mai, K.,
  Wang, Y., \& Liu, Y.
\newblock {ERNIE-RNA}: An {RNA} language model with
  structure-enhanced representations.
\newblock \emph{Nature Communications}, 16:10076, 2025.

\end{thebibliography}


%% ────────────────────────────────────────────
\appendix
\section{RNA family details ($N = 52$)}
\label{sec:appendix_families}

Phase~1 families (12, from original preregistration):

\begin{table}[h]
\centering
\small
\begin{tabular}{@{}llcccc@{}}
\toprule
\textbf{Family} & \textbf{Rfam} & \textbf{Stem} & \textbf{Loop} &
  \textbf{RNA-FM ratio} & \textbf{$>$ nuc null} \\
\midrule
tRNA-Phe (yeast)    & RF00005 & 42 & 34 & 2.758 & No \\
tRNA-Ala (human)    & RF00005 & 42 & 34 & 3.208 & Yes \\
5S rRNA             & RF00001 & 72 & 48 & 1.158 & No \\
Hammerhead          & RF00163 & 16 & 32 & 1.328 & No \\
SAM riboswitch      & RF00162 & 52 & 42 & 2.535 & No \\
TPP riboswitch      & RF00059 & 40 & 40 & 1.442 & No \\
SRP RNA             & RF00017 & 44 & 33 & 1.261 & No \\
miR-21              & RF00077 & 64 &  8 & 1.299 & No \\
IRES HCV dom.~II    & RF00061 & 50 & 27 & 1.710 & No \\
HDV ribozyme        & RF00094 & 48 & 28 & 1.135 & No \\
RNase~P             & RF00010 & 42 & 36 & 1.513 & No \\
U2 snRNA            & RF00004 & 46 & 30 & 2.200 & No \\
\bottomrule
\end{tabular}
\end{table}

Expanded families (40 additional):

\begin{table}[h]
\centering
\small
\begin{tabular}{@{}lllcccc@{}}
\toprule
\textbf{Family} & \textbf{Rfam} & \textbf{Class} & \textbf{Stem} & \textbf{Loop} &
  \textbf{RNA-FM ratio} & \textbf{$>$ nuc null} \\
\midrule
6S RNA              & RF00013 & ncRNA       & 42 & 34 & 1.823 & No \\
7SK RNA             & RF00100 & ncRNA       & 38 & 31 & 1.593 & No \\
Bacterial SRP       & RF00169 & ncRNA       & 43 & 22 & 2.179 & No \\
c-di-GMP ribosw.    & RF01051 & Riboswitch  & 36 & 27 & 1.826 & No \\
Cobalamin ribosw.   & RF00174 & Riboswitch  & 38 & 36 & 2.214 & No \\
Corona 5' UTR       & RF03120 & Cis-reg     & 51 & 33 & 1.252 & No \\
Corona s2m          & RF00164 & Cis-reg     & 56 & 27 & 1.268 & No \\
CRISPR leader       & RF01316 & CRISPR      & 25 & 30 & 1.121 & No \\
CrPV IRES           & RF00458 & IRES        & 40 & 24 & 1.575 & No \\
Fluoride ribosw.    & RF01734 & Riboswitch  & 36 & 28 & 1.473 & No \\
FMN riboswitch      & RF00050 & Riboswitch  & 52 & 26 & 1.889 & No \\
glmS ribozyme       & RF00234 & Ribozyme    & 44 & 38 & 1.252 & No \\
Glycine ribosw.     & RF00504 & Riboswitch  & 26 & 32 & 1.167 & No \\
Group~I intron P4P6 & RF00028 & Intron      & 46 & 31 & 1.827 & No \\
Group~II intron D5  & RF00029 & Intron      & 38 & 21 & 1.011 & No \\
Hatchet ribozyme    & RF02678 & Ribozyme    & 26 & 31 & 1.865 & No \\
HCV IRES III        & RF00061 & IRES        & 40 & 24 & 1.480 & No \\
Histone 3' stem     & RF00032 & Cis-reg     & 10 & 10 & 4.078 & No \\
IRE stem-loop       & RF00037 & Cis-reg     & 21 &  9 & 2.006 & No \\
Lysine ribosw.      & RF00168 & Riboswitch  & 40 & 38 & 1.581 & No \\
Manganese ribosw.   & RF01786 & Riboswitch  & 40 & 24 & 1.008 & No \\
miR-122             & RF00246 & miRNA       & 44 & 28 & 1.252 & No \\
miR-155             & RF00253 & miRNA       & 44 & 22 & 2.529 & No \\
miR-let-7           & RF00027 & miRNA       & 48 & 27 & 0.952 & No \\
Pistol ribozyme     & RF02679 & Ribozyme    & 28 & 38 & 1.673 & No \\
preQ1 ribosw.       & RF00522 & Riboswitch  & 26 & 25 & 2.757 & No \\
Purine ribosw.      & RF00167 & Riboswitch  & 36 & 21 & 2.353 & No \\
SAH riboswitch      & RF01057 & Riboswitch  & 32 & 28 & 1.252 & No \\
SECIS element       & RF00031 & Cis-reg     & 30 & 28 & 1.166 & No \\
T-box leader        & RF00230 & Cis-reg     & 38 & 26 & 2.651 & No \\
THF riboswitch      & RF01831 & Riboswitch  & 52 & 30 & 1.150 & No \\
tmRNA               & RF00023 & ncRNA       & 56 & 29 & 2.682 & No \\
Twister ribozyme    & RF02681 & Ribozyme    & 30 & 37 & 1.584 & No \\
U1 snRNA            & RF00003 & snRNA       & 44 & 31 & 1.721 & No \\
U4 snRNA            & RF00015 & snRNA       & 38 & 34 & 2.101 & No \\
U5 snRNA            & RF00020 & snRNA       & 44 & 18 & 2.864 & No \\
U6 snRNA            & RF00026 & snRNA       & 40 & 30 & 1.091 & No \\
Vault RNA           & RF00006 & ncRNA       & 44 & 29 & 1.264 & No \\
Y RNA               & RF00019 & ncRNA       & 36 & 32 & 1.979 & No \\
ZMP riboswitch      & RF01750 & Riboswitch  & 30 & 26 & 1.961 & No \\
\bottomrule
\end{tabular}
\end{table}


\section{HTT CAG-repeat full results}
\label{sec:appendix_htt}

HTT exon~1 constructs from NM\_002111.7: 30\,nt 5$'$ flank
(from ATG start codon), variable CAG repeat, 154\,nt 3$'$ flank
(including CCG repeat and downstream coding sequence).  Structures
predicted by ViennaRNA RNAfold.

\begin{table}[h]
\centering
\small
\caption{HTT construct properties.}
\begin{tabular}{@{}lcccc@{}}
\toprule
\textbf{Variant} & \textbf{Label} & \textbf{Length (nt)} &
  \textbf{MFE (kcal/mol)} & \textbf{Paired positions} \\
\midrule
CAG17 & Normal       & 235 & $-$102.3 & 148 \\
CAG21 & Normal       & 247 & $-$105.5 & 158 \\
CAG36 & Pathogenic   & 292 & $-$124.0 & 184 \\
CAG40 & Pathogenic   & 304 & $-$129.0 & 192 \\
CAG60 & Pathogenic   & 364 & $-$154.0 & 232 \\
\bottomrule
\end{tabular}
\end{table}

\begin{table}[h]
\centering
\small
\caption{Mutation sensitivity ratio per model per variant.
  Bold entries exceed the nucleotide-stratified composition null
  (95th percentile, 200 permutations).}
\begin{tabular}{@{}lccccc@{}}
\toprule
\textbf{Model} & \textbf{CAG17} & \textbf{CAG21} &
  \textbf{CAG36} & \textbf{CAG40} & \textbf{CAG60} \\
\midrule
RNA-FM     & \textbf{1.753} & 1.928 & 1.982 & \textbf{1.951} & 1.355 \\
RiNALMo    & 1.032 & 1.033 & 1.025 & 1.028 & 1.036 \\
ERNIE-RNA  & 1.077 & 1.088 & 1.079 & 1.095 & 1.089 \\
SpliceBERT & 0.996 & 0.997 & 0.991 & 0.996 & 1.001 \\
UTR-LM     & 1.121 & 1.144 & 1.157 & 1.164 & \textbf{1.167} \\
\bottomrule
\end{tabular}
\end{table}

\begin{table}[h]
\centering
\small
\caption{Composition null 95th percentiles per model per variant.}
\begin{tabular}{@{}lccccc@{}}
\toprule
\textbf{Model} & \textbf{CAG17} & \textbf{CAG21} &
  \textbf{CAG36} & \textbf{CAG40} & \textbf{CAG60} \\
\midrule
RNA-FM     & 1.723 & 1.950 & 2.215 & 1.712 & 1.778 \\
RiNALMo    & 1.087 & 1.092 & 1.089 & 1.098 & 1.099 \\
ERNIE-RNA  & 1.085 & 1.093 & 1.128 & 1.143 & 1.195 \\
SpliceBERT & 1.054 & 1.042 & 1.058 & 1.059 & 1.072 \\
UTR-LM     & 1.179 & 1.190 & 1.177 & 1.168 & 1.153 \\
\bottomrule
\end{tabular}
\end{table}

\begin{table}[h]
\centering
\small
\caption{Attention-contact correlation (Spearman $\rho$, best
  head/layer) per variant.}
\begin{tabular}{@{}lccccc@{}}
\toprule
\textbf{Model} & \textbf{CAG17} & \textbf{CAG21} &
  \textbf{CAG36} & \textbf{CAG40} & \textbf{CAG60} \\
\midrule
RNA-FM     & 0.041 & 0.037 & 0.034 & 0.034 & 0.033 \\
RiNALMo    & 0.059 & 0.056 & 0.054 & 0.053 & 0.049 \\
NTv2       & 0.106 & 0.092 & 0.181 & 0.171 & 0.150 \\
ERNIE-RNA  & 0.077 & 0.074 & 0.069 & 0.068 & 0.062 \\
SpliceBERT & 0.035 & 0.030 & 0.039 & 0.039 & 0.039 \\
UTR-LM     & 0.038 & 0.037 & 0.041 & 0.041 & 0.039 \\
\bottomrule
\end{tabular}
\end{table}

\begin{table}[h]
\centering
\small
\caption{Embedding distance from wild-type (CAG17) reference.
  Cosine distance of full-sequence mean embedding.}
\begin{tabular}{@{}lcccc@{}}
\toprule
\textbf{Model} & \textbf{CAG21} & \textbf{CAG36} &
  \textbf{CAG40} & \textbf{CAG60} \\
\midrule
RNA-FM     & 0.0000 & 0.0000 & 0.0000 & 0.0000 \\
RiNALMo    & 0.0016 & 0.0217 & 0.0294 & 0.0679 \\
NTv2       & 0.0002 & 0.0061 & 0.0070 & 0.0127 \\
ERNIE-RNA  & 0.0006 & 0.0072 & 0.0091 & 0.0229 \\
SpliceBERT & 0.0014 & 0.0128 & 0.0137 & 0.0424 \\
UTR-LM     & 0.0003 & 0.0050 & 0.0064 & 0.0158 \\
\bottomrule
\end{tabular}
\end{table}


\end{document}
